% Proposed B.122. DE-owned isolated source; central integration belongs to root.
\subsection*{The completed centralizer in the actual intermediate amalgam}

\statement{B.122}{Proved}{The selected-cone centralizer in the full completion of the intermediate boundary amalgam is exactly the closure of its actual discrete source centralizer}
\label{stmt:B122}
Use the complete marked intermediate group of B.120,
\[
 M=R*_{H}W_{\partial},\qquad
 H=C_2(c)*\mathbb Z^2(k,T),\qquad c=aCa^{-1},
\]
with the actual selected regional source $L\cong F_5$ and
\[
 \begin{gathered}
 W_{\partial}=\langle D,u,k,j\mid D^2=j^2=1,
 [k,D]=[k,u]=1,\ juj=u^{-1}\rangle,\\
 c=jDj,\qquad T=uD,\qquad k'=jkj.
 \end{gathered}
\]
All completions below are full profinite completions embedded by
the proved induced-topology comparisons. Closures are taken in
$\widehat M$. Then
\begin{equation}\label{eq:B122-closure-centralizers}
 \begin{aligned}
 C_{\widehat M}(c)
   &=\overline{C_M(c)}
     =\overline{\langle L,k',c\rangle},\\
 C_{\widehat M}(D)
   &=\overline{C_M(D)}
     =\overline{\langle jLj,k,D\rangle}.
 \end{aligned}
\end{equation}
Thus the completed centralizer contains no elements beyond the
closure of the actual discrete centralizer from B.120. This is a
closure equality. It does not assert that the full completion of
$F_6\times C_2$ embeds under its discrete centralizer map, or that
the completed centralizer is a free profinite amalgam. The four
remaining prefix relations defining the actual final group are
not imposed in this statement.

\begin{proof}
\emph{B.122.1. The corner decomposition has a proper full completion.}
The isolated corner is the actual ordinary amalgam
\begin{equation}\label{eq:B122-corner-split}
 W_{\partial}=A*_{\langle u\rangle}V,\qquad
 A=(C_2(D)*\mathbb Z(u))\times\mathbb Z(k),\qquad
 V=\mathbb Z(u)\rtimes C_2(j),
\end{equation}
where $j$ acts on $u$ by inversion. Both edge maps are injective.
We give the finite-target comparison explicitly. Prescribe any
finite targets of $A$ and $V$, and let the orders of the two
images of $u$ be $n_A,n_V$. Put $n=\operatorname{lcm}(n_A,n_V)$.
Refine the first target by the map
\[
 A\longrightarrow C_n,\qquad D,k\longmapsto1,\quad
 u\longmapsto\text{a generator},
\]
and the second by the finite dihedral quotient
$V\to C_n\rtimes C_2$. In both refined finite vertex images,
the actual image of $u$ now has exactly order $n$. Their edge
kernels are therefore the identical subgroup $\langle u^n\rangle$.

For completeness, two finite groups with a common embedded subgroup
have a finite permutation image of their amalgam which is injective
on both vertices. Choose free actions of the two finite groups on
one set of cardinality divisible by their orders. The restrictions
to the common subgroup are unions of equally many regular actions.
Conjugating one permutation representation aligns those restrictions.
The resulting representations agree on the edge and remain faithful
on the two finite vertices. Applying this construction to the refined
targets proves that each vertex of \eqref{eq:B122-corner-split}
has its full induced topology in $W_{\partial}$. The cyclic edge
already has full topology in each vertex, by its retraction in $A$
and its index-two inclusion in $V$. The universal property on
finite targets consequently gives the proper profinite amalgam
\begin{equation}\label{eq:B122-completed-corner}
 \widehat{W_{\partial}}=
 \widehat A\amalg_{\widehat{\mathbb Z}(u)}\widehat V.
\end{equation}
No residual-finiteness theorem for a general amalgam is needed.

The factor completions are
\[
 \widehat A=(C_2(D)\amalg\widehat{\mathbb Z}(u))
                         \times\widehat{\mathbb Z}(k),
 \qquad
 \widehat V=\widehat{\mathbb Z}(u)\rtimes C_2(j).
\]
The direct-product comparison follows by restricting any finite
target to its commuting factors and taking the product of their
restriction kernels. The dihedral comparison follows from the
index-two cyclic subgroup. In particular every edge stabilizer in
the standard profinite tree of \eqref{eq:B122-completed-corner}
is a conjugate of the torsion-free group
$\widehat{\mathbb Z}(u)$.

\emph{B.122.2. The exact completed vertex centralizers.}
We use the standard profinite tree for a proper finite graph of
profinite groups, as constructed by
\citet[Proposition~3.8]{B_ZalesskiiMelnikov1989}. Their
Theorem~2.8 states that a nonempty fixed subgraph is connected;
only its order-two case is needed here. The original proof on
printed pp.~413--414 proves this case through the mod-two chain
complex. Connectedness here is connectedness of profinite graphs,
meaning that all finite quotient graphs are connected, not
connectedness of the underlying totally disconnected space.

In the corner tree, $D$ fixes the base $\widehat A$-vertex and
cannot fix an edge, because the edge stabilizers are torsion free.
Its fixed subgraph is connected and has no edges, and therefore
is a single vertex. Every element centralizing $D$ preserves that
unique vertex, so it lies in $\widehat A$. Herfort and Ribes'
finite-factor centralizer theorem gives
\begin{equation}\label{eq:B122-corner-centralizers}
 \begin{aligned}
 C_{\widehat{W_{\partial}}}(D)
   &=\langle D\rangle\times\widehat{\mathbb Z}(k),\\
 C_{\widehat{W_{\partial}}}(c)
   &=\langle c\rangle\times\widehat{\mathbb Z}(k')
     =\overline{C_{W_{\partial}}(c)}.
 \end{aligned}
\end{equation}
Here the first equality uses
\citet[Theorem~B$'$]{A_HR1985} in
$C_2(D)\amalg\widehat{\mathbb Z}(u)$, and the second is its
conjugate by the actual $j$.

B.112, with the entire selected-gap conjugation, already gives
\begin{equation}\label{eq:B122-regional-centralizer}
 C_{\widehat R}(c)=\widehat L\times\langle c\rangle
                 =\overline{C_R(c)}.
\end{equation}
Also, by B.118,
\[
 \widehat H=C_2(c)\amalg\widehat{\mathbb Z}^{\,2}(k,T).
\]
The second factor is torsion free. Theorems~A$'$ and B$'$ of
\citet{A_HR1985} therefore show that every involution of
$\widehat H$ is $\widehat H$-conjugate to $c$, and
\begin{equation}\label{eq:B122-edge-centralizer}
 C_{\widehat H}(c)=\langle c\rangle.
\end{equation}
These assertions concern the actual embedded full edge completion.

\emph{B.122.3. The compact fixed graph and its one-edge quotient.}
By B.120 the full completion of $M$ is the proper amalgam
\[
 \widehat M=\widehat R\amalg_{\widehat H}
                         \widehat{W_{\partial}}.
\]
Both edge maps are proper. On the corner side mirror parity
is zero on $H$ and one on $j$. On the regional side the
actual character to $C_2$ sending $a,b$ to one and
$C,k,h,d,e$ to zero is zero on $c,k,T$ and nonzero on $a$.
These finite characters also witness properness after completion.
Let $\mathscr T$ be its standard profinite tree. Its two vertex
spaces are $\widehat M/\widehat R$ and
$\widehat M/\widehat{W_{\partial}}$, and its edge space is
$\widehat M/\widehat H$. In particular the edge space is compact.
The actual base edge is fixed by $c$. Thus
\[
 \Sigma=\mathscr T^{\langle c\rangle}
\]
is nonempty and connected by the fixed-subgraph theorem just cited.
Its edge space is a closed subset of the compact edge space of
$\mathscr T$, hence is compact as well.

Every vertex of $\Sigma$ is incident with a fixed edge. Indeed,
the union of the two endpoint images of its compact edge space is
closed. If a vertex were outside it, a nonempty clopen vertex
neighborhood disjoint from all endpoints could be separated from
the nonempty endpoint set. Sending that neighborhood to one point
and the rest of the graph to another would give a disconnected
finite quotient graph, contradicting connectedness. This is also
exactly the compact-edge lemma in
\citet[Lemma~2.14]{B_ZalesskiiMelnikov1989}. The compactness
hypothesis is essential; no finite-path argument is being used.

Put $Z=C_{\widehat M}(c)$. This compact group acts continuously
on $\Sigma$. At any fixed vertex choose an incident fixed edge.
Its stabilizer is a conjugate $H_e$ of $\widehat H$ containing
$c$. Every other incident fixed edge has the form $g e$, with
$g$ in that vertex stabilizer and $g^{-1}cg\in H_e$. By the
unique involution class in $H_e$, choose $t\in H_e$ with
\[
 g^{-1}cg=tct^{-1}.
\]
Then $gt$ centralizes $c$ and sends $e$ to $ge$. Hence the
$Z$-stabilizer of each vertex is transitive on its incident fixed
edges.

The profinite orbit graph $Z\backslash\Sigma$ is connected.
The two vertex types remain distinct. Every vertex has exactly
one incident edge in this quotient. Thus its two endpoint maps
from the edge space to the respective vertex-type spaces are
continuous bijections, hence homeomorphisms. If its edge space
had more than one point, a nontrivial clopen partition of that
space, transported by both endpoint homeomorphisms, would give
a disconnected finite quotient graph. Connectedness excludes
this. Consequently
\begin{equation}\label{eq:B122-one-edge-orbit}
 Z\backslash\Sigma\quad\text{is exactly one edge and its two vertices}.
\end{equation}

\emph{B.122.4. Closed generation by the two actual vertex centralizers.}
At the base edge write
\[
 A_c=C_{\widehat R}(c),\qquad
 B_c=C_{\widehat{W_{\partial}}}(c),\qquad
 E_c=C_{\widehat H}(c)=\langle c\rangle.
\]
By \eqref{eq:B122-one-edge-orbit}, the fixed graph is the actual
coset graph with vertex spaces $Z/A_c$, $Z/B_c$ and edge space
$Z/E_c$. These identifications are homeomorphisms: the orbit
maps are continuous bijections from compact spaces to Hausdorff
spaces.

Let $Z_0=\overline{\langle A_c,B_c\rangle}$ in $Z$.
There are well-defined continuous maps
\[
 Z/A_c\longrightarrow Z/Z_0,\qquad
 Z/B_c\longrightarrow Z/Z_0,\qquad
 Z/E_c\longrightarrow Z/Z_0,
\]
all sending a coset represented by $z$ to $zZ_0$. The two
endpoints of every edge have the same image. They therefore
define a continuous graph map from $\Sigma$ onto the edgeless
profinite graph $Z/Z_0$, with edges contracted to vertices.
Every finite quotient of a connected profinite graph is connected.
An edgeless profinite graph with this property is a singleton.
Thus $Z/Z_0$ is a singleton and
\begin{equation}\label{eq:B122-generation}
 C_{\widehat M}(c)=\overline{\langle
                 C_{\widehat R}(c),
                 C_{\widehat{W_{\partial}}}(c)\rangle}.
\end{equation}
This uses only connectedness and compact coset spaces. It does
not identify the fundamental group of the fixed graph.

Substitute \eqref{eq:B122-corner-centralizers} and
\eqref{eq:B122-regional-centralizer} into
\eqref{eq:B122-generation}. B.120's exact discrete formula
$C_M(c)=\langle C_R(c),C_{W_{\partial}}(c)\rangle$
then gives the first line of
\eqref{eq:B122-closure-centralizers}. Conjugating by the actual
$j$ gives its second line. The proof establishes the closed image
of the discrete centralizer, without assuming that this discrete
centralizer has its full induced topology. It makes no assertion
after the four remaining prefix relations are imposed.
\end{proof}
